\section{The model and its exact transform}\label{spm:sec:model}
A matrix is \emph{packed} if it has no zero row or column. Let $\mathcal A_n$ be the set of symmetric packed matrices with nonnegative integer entries and total entry sum $n$. The dimension is allowed to vary. Rows and columns are ordered, and simultaneous permutations are not identified. The empty matrix is the sole object of weight zero. For $n>0$, the dimension is at most $n$, so the set is finite. An off-diagonal entry of value $m$ contributes $2m$ to the total sum, while a diagonal entry contributes its value once.

Put $A_n=|\mathcal A_n|$. The first values are
\[
1,\ 1,\ 3,\ 9,\ 33,\ 125,\ 531,\ 2349,\ 11205,\ 55589,\ 291423.
\]
This is OEIS A138178~\cite{oeis138178}. The entry also identifies these objects with normal semistandard Young tableaux, whose used alphabet is an initial interval. The matrix model, rather than any quotient by relabeling, is used throughout this report.

For $M\in\mathcal A_n$, let $D(M)$ be its actual dimension and $T(M)$ its diagonal sum, and define
\[
 A_n(u,v)=\sum_{M\in\mathcal A_n}u^{D(M)}v^{T(M)},\qquad A_0(u,v)=1.
\]
The auxiliary quantities
\begin{equation}\label{spm:eq:notation}
 q=\frac{u}{1+u},\qquad L=L(u)=\log(1+1/u),\qquad
 I_n(v)=n![x^n]e^{vx+x^2/2}
\end{equation}
use the branch of the logarithm positive at $u=1$. In particular $L(1)=\log2$, and $I_n(1)$ is the involution number. Define the coefficient polynomial
\begin{equation}\label{spm:eq:fn}
 f_n(k,v)=[z^n](1-vz)^{-k}(1-z^2)^{-(k^2-k)/2}.
\end{equation}
For integer $k\geq0$, this counts all symmetric $k\times k$ matrices, with zero lines allowed, marking their diagonal sum. For every fixed $n>0$, it is a polynomial in $k$ of degree at most $n$ with zero constant term.

\begin{proposition}[Exact geometric transform]\label{spm:prop:transform}
For $u,v$ in a complex neighborhood of $(1,1)$,
\begin{equation}\label{spm:eq:transform}
 A_n(u,v)=\frac1{1+u}\sum_{k\geq0}q^k f_n(k,v).
\end{equation}
The series converges absolutely for fixed $n$. The marker $u$ records the actual packed dimension, not the auxiliary index $k$.
\end{proposition}
\begin{proof}
Deleting simultaneous zero rows and columns from a $k\times k$ matrix leaves a unique packed matrix. Each particular packed matrix of dimension $d$ has $\binom{k}{d}$ preimages. Since $1-q=(1+u)^{-1}$,
\[
 (1-q)\sum_{k\geq d}q^k\binom{k}{d}
 =\left(\frac{q}{1-q}\right)^d=u^d.
\]
The diagonal sum is unchanged by insertion of zero lines. Summing over the finite family of packed matrices proves the identity. The polynomial degree assertion follows by expanding the logarithm of \eqref{spm:eq:fn}: every atom $k^a z^b$ has $a\leq b$.
\end{proof}

The unmarked geometric generating-function transform is already recorded in OEIS, credited to Vladeta Jovovic in December 2009~\cite{oeis138178}. Its use here, and the short marked extension above, must not be mistaken for a newly discovered unmarked identity.

\subsection{Prior results and the boundary of the model}
Cameron, Prellberg, and Stark~\cite{cps}, Proposition 4.2, prove the leading equivalent for the symmetric \emph{binary} companion $B_n$, OEIS A135588~\cite{oeis135588}:
\begin{equation}\label{spm:eq:priorbinary}
 B_n\sim e^{-L-L^2/4}\frac{I_n(1)P_n}{n!},\qquad L=\log2,
\end{equation}
where $P_n$ is the ordered Bell number. Their Section 2 also obtains a Poisson collision limit in a nonsymmetric incidence-matrix argument based on pairs of random preorders. The collision mechanism therefore has substantial precedent. The present treatment concerns the nonnegative symmetric model, uniform marked expansions to every fixed order, and quantitative refinements of its joint collision law.

Rectangular packed matrices, matrices with decreasing margins, and matrices considered up to row or column permutations are different counting models. Related rectangular and Burge-polynomial frameworks appear in \cite{matrixcomp,caylerian}. Their familiar names and related generating functions do not identify their sequences with $A_n$. A bounded check of relevant public sources and existing reports found no duplicate of the result assembled here; it is not a certificate of worldwide novelty. The source notes accompanying this report describe the public attributions and the limits of that check.

\begin{writenote}[Added 6 October 2026, batch~108: the rectangular counterpart]
The rectangular counterpart of this model on the diagonal of total weight
$2N$ with $N$ rows (OEIS A261784; bundle Report~169, which the source notes
name among the reports they read) is Part~IV (Sections~34--44, labels
\file{mxc:pm:}) of the collection's report
\file{a261781-matrix-compositions}. It uses the devices of this report on a
different model: a geometric weighting of dummy columns that erases zero
lines, the pole of the ordered Bell generating function at $\log(1+1/u)$,
an expansion to every fixed order uniform for complex markers of the
columns and of the cells with entry at least~$2$ (its
Theorem~\file{mxc:pm:thm:marked}, on $|\omega-1|\le2.1$, $|u-1|\le0.02$),
and a joint limit in which the number of such cells is Poisson of mean
$r\log2=1.1046\ldots$, $r=2+W_0(-2e^{-2})$, independent of the Gaussian
column count (Theorem~\file{mxc:pm:thm:joint}). There the expansion is in
powers of $N^{-1}$; here it is in powers of $n^{-1/2}$ around an
involution saddle, and the repeated cells split into diagonal and
off-diagonal ones with means $\log2$ and $(\log2)^2/2$
(Theorem~\ref{spm:thm:joint}). No theorem is shared, and neither report
uses the other. The collection's report
\file{a260700-parabolic-double-cosets} extracts the same ordered Bell pole
at $\log2$ for a different sequence, and the exact pole-lattice formula
for the ordered Bell numbers is Theorem~\file{q2:thm:fubini} of the
transseries volume
\path{Analysis/Transseries/docs/series-and-transseries/Transseries_And_Inversion/transseries_and_inversion.tex}:
the last sentence of Corollary~\ref{spm:cor:unmarked} is its one-pole
truncation, and~\eqref{spm:eq:poles} at $L=\log2$, $j=n$, multiplied
by~$\tfrac12$, is its first equality.
\end{writenote}

\begin{remark}[{[write]} The OEIS entries and the data, 6~October 2026]\label{spm:rem:oeis}
The live entry A138178 \cite{oeis138178} (revision~\#39 of 18~December
2022, the revision the source inspected; read 6~October 2026) is named,
verbatim, ``Number of symmetric matrices with nonnegative integer entries
and without zero rows or columns such that sum of all entries is equal to
n.'' Its offset is~$0$; its 25 data terms $A_0,\ldots,A_{24}$ equal the
source's transcription (the delivered receipt
\file{data/count_receipt.json}), the first eleven being those displayed
above. The entry has two generating functions: an unattributed
inclusion--exclusion form, which is the generating function
of~\eqref{spm:eq:deletion} at $u=v=1$, and Jovovic's
$\sum_{k\ge0}2^{-k-1}(1-x)^{-k}(1-x^2)^{-\binom k2}$ of 9~December 2009,
which is~\eqref{spm:eq:transform} at $u=v=1$ (then $q=\tfrac12$). Its
comments give Wiseman's normal semistandard tableaux (2018) and
examples; it has no asymptotic statement, so Corollary~\ref{spm:cor:unmarked}
is the only asymptotic for A138178 in the entry or in the repository.
Heinz's b-file lists $A_0,\ldots,A_{500}$. The source's retrieval failed
(\file{SOURCES.md}), and it claims no comparison
(Section~\ref{spm:sec:computation}); the write fetched it on
6~October 2026 and computed $A_0,\ldots,A_{500}$ exactly by an
implementation of its own (the recurrence~\eqref{spm:eq:exactrecurrence}
at $v=1$ and $\tfrac12\sum_{k\ge0}2^{-k}k^j=P_j$ for the ordered Bell
numbers): all 501 terms agree. The binary companion A135588
\cite{oeis135588} (revision~\#17 of 22~January 2024) has 26 data terms
$B_0,\ldots,B_{25}$; an inclusion--exclusion count of the write's
(binomial cell factors, as in Section~\ref{spm:sec:computation}) gives
all 26. It has no asymptotic statement either. Nothing in either entry is
corrected, and nothing was submitted to the OEIS. The leading formula and
its corrections are compared with the exact counts in
Remark~\ref{spm:rem:third}.

\emph{The binary precedent.} The write read Cameron--Prellberg--Stark
\cite{cps} (arXiv version~2 of 4~April 2006). Their
Proposition~4.2 (printed p.~14, proof pp.~14--15) states $S_{11}(n)\sim C_sI(n)P(n)/n!$ with
$C_s=\tfrac12e^{-(\log2)^2/4}\approx0.44341$, where $S_{11}(n)$ counts
the symmetric incidence matrices with $n$ ones (A135588, here $B_n$),
$I(n)$ the involutions and $P(n)$ the preorders, that is the ordered Bell
numbers. Since $\tfrac12=e^{-\log2}$, this is~\eqref{spm:eq:priorbinary}.
Their Section~2 (the first proof of their Theorem~2.1, pp.~6--9) shows by
the method of moments that the number $W$ of pairs of points lying in a
common block of both of two independent uniform random preorders converges
to a Poisson$((\log2)^2/2)$ law, so that $\PP(W=0)\to e^{-(\log2)^2/2}$
(their~(8), p.~9):
the nonsymmetric collision precedent the source credits. The source's
account of both is accurate.
\end{remark}

\subsection{Provenance, sources and reading conventions}\label{spm:sec:provenance}
\begin{writenote}[Added 6 October 2026, batch~108]
\emph{Provenance.} This report is Report~238 of the session bundle that
ProveIt's \file{docs/incoming} intake processed as batch~108: the archive
\file{Report238.zip} (604{,}411 bytes, SHA-256
\file{255ba5ee21d1...5f6eac3b}, 30 files in a wrapper directory
\file{Report238/}), arrival commit \file{60f54ea06}, placed in
\file{602e5bd0f}. The text is the delivered \file{article.tex} with its
twelve files \file{sections/*.tex} (952 lines in all, dated 5~October
2026), shipped under the same names. The delivered 22-page PDF and the
checksum manifest \file{MANIFEST.sha256} (29 entries, verified at
placement and again by the write) are not shipped, and the delivered
README is replaced by the report README; all three are retrievable from
the arrival commit. The package records no ProveIt commit; its
\file{SOURCES.md} describes a bounded look at existing reports, naming
the matrix-composition report and bundle Reports~169 and~177 (now Part~IV
of \file{a261781-matrix-compositions} and the batch-109 report
\file{a262810-diagonal-alignments}). The phrase ``private review'' does
not occur in it; the PDF author field is ``Report 238''. It is a
single-source report, so no merge choice was made. The write prefixed
the 110 delivered labels with \file{spm:} and updated the 93 references
to them; it added this subsection, the notes marked [write],
Remarks~\ref{spm:rem:oeis}, \ref{spm:rem:third}
and~\ref{spm:rem:transseries} and Section~\ref{spm:sec:further}. Apart
from the label prefixes no theorem, proof or number of the source was
changed: the added statements are the last of their sections, the added
section comes after the last numbered one, and the added displays are
unnumbered, so every theorem, section and equation keeps its delivered
number.

\emph{The sources, as the write read them} (6~October 2026).
\begin{itemize}\raggedright
\item OEIS A138178 (revision~\#39), its b-file, and A135588
 (revision~\#17) (Remark~\ref{spm:rem:oeis}).
\item Cameron--Prellberg--Stark \cite{cps}, arXiv:math/0510155v2:
 Proposition~4.2 and its proof (pp.~14--15) and the first proof of their
 Theorem~2.1 in Section~2 (pp.~6--9), as in Remark~\ref{spm:rem:oeis}.
\item The transseries volume
 \path{Analysis/Transseries/docs/series-and-transseries/Transseries_And_Inversion/transseries_and_inversion.tex}:
 \file{p0:def:core}, \file{p0:prop:factorial-core},
 \file{p0:thm:core-reversion} with \file{p0:rem:core-instances},
 \file{p0:def:three-inverses}, \file{p0:thm:staircase},
 \file{plt:thm:lw-template} and \file{q2:thm:fubini}, read for
 Remark~\ref{spm:rem:transseries} and the note above.
\item Part~IV of \file{a261781-matrix-compositions}, at the statements
 named in the note above.
\item Not read by the write: Munarini--Poneti--Rinaldi \cite{matrixcomp}
 and Cerbai--Claesson \cite{caylerian}. What is said about them (related
 models, not the transpose-symmetric count) is the source's.
\end{itemize}

\emph{What was checked at intake and by the write.} The batch-108 dossier
read the manuscript and found no error; it checked by hand the
total-variation constant of Theorem~\ref{spm:thm:TV} from the first mean
corrections, and reran the package on a copy (Windows, Python~3.14.4,
SymPy~1.14.0, mpmath~1.3.0): \file{build.py} with \verb|--verify-only|
verified 29 files; \file{exact_counts.py} (to $n=640$),
\file{symbolic_coefficients.py} and \file{diagnostics.py} reproduced the
three delivered receipts; \file{guard_tests.py} stops on Windows at its
own path guard; the PDF and ZIP rebuilds were not run. The write reran the
three mathematical programs from the shipped files (README, Route~B):
their outputs equal the receipts byte for byte. It re-derived by hand,
and in part with SymPy, $C_1$ and $\mathcal B_1$, $G_1$, $H_1$ of
Section~\ref{spm:sec:allorders}, the agreement of~\eqref{spm:eq:D2} with
the delivered symbolic receipt, the cumulants of
Proposition~\ref{spm:prop:cumulants} and the residue~\eqref{spm:eq:trace-residue},
the coefficients~\eqref{spm:eq:Sstar}, \eqref{spm:eq:C1star}
and~\eqref{spm:eq:binaryS}--\eqref{spm:eq:binaryprob}, the mass
correction~\eqref{spm:eq:masscorrection} and the constant
of~\eqref{spm:eq:sharpTV}, the constants~\eqref{spm:eq:cbd}, and the
coefficients~\eqref{spm:eq:alpha}--\eqref{spm:eq:gamma} (also by the
triangular recurrence of the transseries volume,
Remark~\ref{spm:rem:transseries}); and it confirmed the monotonicity
argument of Section~\ref{spm:sec:inverse}. Finite computations are
observations, not certificates.

\emph{Relation to the repository.} Before batch~108 no file of the
repository mentioned A138178, A135588 or symmetric packed matrices. No
statement of this report is formalized in Lean or Rocq, and its place in
the collection confers no formal status. The ordered Bell numbers used in
Corollary~\ref{spm:cor:unmarked} are defined in Lean as
\file{Fabius.fubini}, with their exponential generating function
(\file{Fabius.two_sub_exp_mul_egfA_fubini}), in
\path{Analysis/FabiusFunction/Lean/FabiusFunction/OrderedBell.lean}; no
asymptotic statement about them is formalized there. The neighbouring
reports are named in the note above and in the report README.
\end{writenote}
\begin{writenote}[What is not claimed, collected from the source]
The unmarked geometric transform is Jovovic's (OEIS, 2009) and ``must not
be mistaken for a newly discovered unmarked identity''; the binary leading
equivalent is Cameron--Prellberg--Stark's, ``recovered here as a
consistency check'', and the Poisson collision mechanism ``has substantial
precedent''. The expansions are local in the dimension and trace markers:
the trace marker is not carried through $v=0$, the order $J$ is fixed
independently of $n$, and ``no uniform estimate in a growing order is
asserted''. The diagnostics at $(u,v)=(2,2)$ and $(1,\tfrac12)$ are ``not
an enlargement'' of the neighbourhood of Theorem~\ref{spm:thm:main}. No
span-one local limit of $T_n$, and no total-variation approximation of
$T_n$ by an unrestricted Poisson law, is asserted. The signed corrections
are ``not claimed positive''; the adjusted Poisson law is not claimed
optimal, and neither the $O(n^{-1})$ constant nor a starting index is
effective. The inverse is ``the inverse of the specified smooth
approximation, not a definition of a noninteger combinatorial count'';
\eqref{spm:eq:rounding} ``does not claim that the ceiling of the bare
displayed truncation is always correct''; Theorem~\ref{spm:thm:ceilings}
gives no numerical constants and ``is an asymptotic bound, not a certified
finite-input inverse program''. The bounded source check ``is not a
certificate of worldwide novelty'', and no global priority, complete
transseries or effective onset is asserted. The exact checks have hard caps
that are implementation boundaries, ``not a range restriction on the
mathematical theorems''; the symbolic checks ``do not constitute a
machine-checked proof''; the decimal outputs are noncertified; and
reproducibility is claimed for the recorded toolchain only.
\end{writenote}
\medskip\noindent\begin{minipage}{\textwidth}
\emph{Reading conventions} (letters the manuscript reuses, with the
tempting false readings; no symbol of the source was renamed; symbols of
the transseries volume carry the subscript ``vol'' in
Remark~\ref{spm:rem:transseries}).\par
\smallskip
\begingroup\small\renewcommand{\arraystretch}{1.1}
\noindent\begin{tabular}{@{}>{\raggedright\arraybackslash}p{0.75in}>{\raggedright\arraybackslash}p{3.45in}>{\raggedright\arraybackslash}p{1.95in}@{}}
\toprule
Symbol & Meaning here (where) & Not to be read as\\
\midrule
$L$ & $\log(1+1/u)$, so $\log2$ at $u=1$ and throughout Sections~\ref{spm:sec:TV}--\ref{spm:sec:inverse} & the volume's target $L_{\rm vol}$\\
$q$ & $u/(1+u)=e^{-L}$ & a Poisson mean\\
$w$ & collision marker (Section~\ref{spm:sec:collision}); Gaussian angular variable $w$, $W$ (Section~\ref{spm:sec:allorders}); $W(2y/(eL^2))$ (Section~\ref{spm:sec:inverse}) & each other\\
$D(M)$, $D_n$ & dimension of a matrix, of a uniform one & $D_2$, the $n^{-1}$ log-coefficient~\eqref{spm:eq:D2}; the exponent $D_J$ in~\eqref{spm:eq:taylorbound}\\
$S(L,v)$, $S_*$; $S(M)$, $S_n$ & amplitude exponents \eqref{spm:eq:S}, \eqref{spm:eq:Sstar}; repeated off-diagonal cells & each other\\
$R$, $R_n$, $R_j$ & pole-separating radius~\eqref{spm:eq:R}; repeated diagonal cells; inverse coefficients~\eqref{spm:eq:inverseall} & each other\\
$K$ & auxiliary integer (Section~\ref{spm:sec:saddles}); compact marker set (Theorem~\ref{spm:thm:collision}) & each other; not the dimension $D$\\
$B$, $B_n$, $B_J$ & correction~\eqref{spm:eq:B}; binary count (A135588); remainder constant (Theorem~\ref{spm:thm:ceilings}) & each other\\
$a$, $b$, $c$, $d$ & in Section~\ref{spm:sec:inverse} the constants~\eqref{spm:eq:cbd}; elsewhere $a=(v^2-1)/2$ or~\eqref{spm:eq:C1star}, $c=(2\zeta-1)/4$, $a=\log u$, $b=\log v$ (Section~\ref{spm:sec:fluctuations}), $b_{nj}$, $d_m$, $d_n(k)$, and generic constants $c$, $C$ & each other; $a_{\rm vol}$, $b_{\rm vol}$, $d_{\rm vol}$\\
$h$, $H$, $H_j$, $H_\pm$ & $n^{-1/2}$; angular exponent~\eqref{spm:eq:H}; Charlier-type corrections~\eqref{spm:eq:Charlier}; envelopes in the proof of Theorem~\ref{spm:thm:ceilings} & each other; $h_{\rm vol}$\\
$Q$, $Q_j$ & $w+1$ in~\eqref{spm:eq:Wseed}; normalized collision polynomials~\eqref{spm:eq:Qpolynomials} & each other; the integer polynomials $Q_n$ of \file{COMPUTATION.md}\\
$N$ & Lambert seed~\eqref{spm:eq:Wseed}; also the bound $N\ge|\zeta-1|$ in~\eqref{spm:eq:collision-envelope} & each other; $N_\ast^{\rm vol}$\\
$M$, $M_n$ & a matrix, the bound $M\ge|w-1|$; the main term~\eqref{spm:eq:gammaexact} & each other\\
$P_n$, $P_j(u)$, $P_1$, $P_2$ & ordered Bell numbers; ordered-Bell transforms (Section~\ref{spm:sec:computation}); Poisson variables & each other\\
$\lambda$, $\mu$, $\kappa$, $\rho$ & Poisson means $L$, $L^2/2$, $\kappa=L+L^2$; $\rho(h)$ of~\eqref{spm:eq:scaling} and the radius $\rho>1$ of Section~\ref{spm:sec:TV}; $\lambda$ is also a scale in Lemma~\ref{spm:lem:gammabound} & each other; $\Lambda$ is the involution law of Section~\ref{spm:sec:saddles}, not $\Lambda_{\rm vol}$\\
$E_n$, $\mathcal E$ & pole error~\eqref{spm:eq:mainpluserror}, generating-function remainder (Section~\ref{spm:sec:TV}); Gaussian functional & each other; $\E$\\
$r$, $s$, $t$ & radius~\eqref{spm:eq:r}, Poisson arguments, gamma variable; also markers in proofs & each other\\
\bottomrule
\end{tabular}
\endgroup
\end{minipage}\par
